\startchapter[title={Relevamiento de los modelos lineales}, reference=Anexo_Relevamiento]

Antes de definir los escenarios de estrés (Sección~\in[sec:escenarios_estres]) se
realizó una serie de estudios exploratorios sobre los modelos lineales, cuyas
conclusiones resume la Sección~\in[sec:relevamiento]. Un test responde si una instancia
determinada se resuelve o no; un estudio, en cambio, recorre una familia de
instancias y muestra a partir de qué tamaño, y por qué, un modelo deja de
resolver en tiempos razonables. En cada estudio se varió un único factor
por vez (tamaño del almacén, cantidad de lotes, pedidos por lote, cantidad
de {\em pickers} y sus velocidades), manteniendo fijo el resto.

% TODO: Tengo que verificar ese porcentaje. No es claro.
En todos los estudios, el almacén es una grilla de $n \times n$ ubicaciones,
con el depósito y el punto de entrega en esquinas opuestas. Cada pedido se
ubica al azar en una de las ubicaciones, con una semilla fija, los pedidos se
reparten entre los lotes por turnos y las fechas de entrega se distribuyen
entre el 60\% y el 140\% de una estimación de la duración total de la
operación, de modo que algunos pedidos terminen antes de su fecha y otros
después.

\startsection[title={Qué dimensión limita al modelo Split}]

El primer estudio partió de una instancia base de 6 lotes, 2 {\em pickers},
26 pedidos por lote, una grilla de 7×7 y velocidades uniformes, y movió un
eje por vez con un tope de 300 segundos por resolución. El
Cuadro~\in[tab:relev_ejes] resume el resultado.

\startplacetable[reference=tab:relev_ejes, location=here, title={Eje que se movió en cada barrido y resultado del modelo Split, con un tope de 300 segundos por resolución.}]
	\bTABLE[split=yes, align=middle, frame=on, width=fit]
	\setupTABLE[row][1][background=color, backgroundcolor=darkgray, style=bold]
	\setupTABLE[column][3][width=7.5cm, align=flushleft]
	\bTR \bTH Eje \eTH \bTH Rango \eTH \bTH Resultado \eTH \eTR
	\bTR \bTD Pedidos por lote \eTD \bTD 10 a 50 \eTD \bTD Resolvió en todo el rango (con 50: Etapa~1 en 46,5~s, Etapa~2 en 14,1~s) \eTD \eTR
	\bTR \bTD Lotes \eTD \bTD 6 a 8 \eTD \bTD Con 8 lotes, la Etapa~2 alcanzó el tope tras explorar 518.628 nodos \eTD \eTR
	\bTR \bTD {\em Pickers} \eTD \bTD 2 a 5 \eTD \bTD Resolvió en todo el rango, aunque no de forma monótona (Etapa~2: 12,9~s con 3, 25,4~s con 4 y 7,3~s con 5) \eTD \eTR
	\bTR \bTD Grilla \eTD \bTD 7×7 a 15×15 \eTD \bTD Con 15×15, la Etapa~1 superó el tope \eTD \eTR
	\eTABLE
\stopplacetable

Solo dos ejes impidieron la resolución, y cada uno lo hizo en una etapa
distinta: la cantidad de lotes aumenta el tiempo de la Etapa~2 (asignación y
secuencia), mientras que el tamaño del almacén el de la Etapa~1
(ruteo). Agregar pedidos a cada lote, en cambio, no tuvo impacto significativo: 
la Etapa~2 exploró 153.934 nodos con 10 pedidos por lote y
4.144 con 26. Esta separación entre las dos etapas orientó el resto de la
exploración, ya que permite elegir qué etapa se pone bajo presión en cada
prueba.

\stopsection

\startsection[title={Exploración de instancias más realistas}]

Los estudios siguientes buscaron que las instancias se parecieran a un
almacén real sin volverse imposibles de resolver.

{\bf Pesos de la grilla.} Con el mismo peso en todos los desplazamientos,
muchas rutas distintas empatan en costo, algo poco representativo de un
almacén, donde recorrer un pasillo no cuesta lo mismo que cruzar entre
pasillos. Con pesos asimétricos (2 y 6), la Etapa~1 llegó a ser 3,7 veces
más rápida en la grilla de 7×7, pero la ventaja desapareció en la de
15×15, por lo que no constituye una mejora de rendimiento. Se adoptaron
pesos levemente distintos según la dirección del desplazamiento (4 y 5),
por su realismo y no por su efecto en los tiempos.

{\bf Velocidades de los {\em pickers}.} Si bien en un almacén real los
operarios tienen velocidades distintas, estas velocidades no son muy distantes,
ya que todos provienen de un mismo perfil y usan el mismo equipamiento. Por este motivo, 
consideramos que un perfil de velocidades realista no tiene una velocidad distinta 
por cada operario, sino unas pocas bandas de velocidad. Este cambio
resultó el más costoso del perfil (Figura~\in[fig:relev_velocidades]).
Con 3 {\em pickers}, pasar de una única velocidad a dos bandas multiplicó
por 7,6 el tiempo de la Etapa~2, y una tercera banda no lo empeoró. Con 4
{\em pickers}, cualquier configuración con más de una velocidad superó el
tope de 300 segundos. En el estudio las bandas se ubicaron entre 0,7 y
1,3; el perfil final adoptó dos bandas, de 0,8 y 1,2.

\placefigure[force][fig:relev_velocidades]{Tiempo de la Etapa~2 según la cantidad de velocidades distintas en la cuadrilla (6 lotes, 26 pedidos por lote, grilla de 7×7). Las barras sin relleno no terminaron dentro del tope de 300 segundos.}
{\startMPcode
numeric w, h, u, bw; w := 11cm; h := 6cm; u := h/320; bw := 1.0cm;
drawarrow (0,0)--(w+0.4cm,0) withpen pencircle scaled 1pt;
drawarrow (0,0)--(0,h+0.4cm) withpen pencircle scaled 1pt;
for t = 0, 100, 200, 300:
  draw (0,t*u)--(-0.12cm,t*u) withpen pencircle scaled .5pt;
  label.lft(textext("\tfx " & decimal t), (-0.12cm,t*u));
endfor
draw (0,300u)--(w,300u) dashed evenly withpen pencircle scaled .5pt withcolor .5white;
label.urt(textext("\tfx tope: 300 s"), (0.1cm,300u));
def barra(expr x, y, c) =
  fill (x,0)--(x+bw,0)--(x+bw,y*u)--(x,y*u)--cycle withcolor c;
enddef;
def barracortada(expr x, y, c) =
  draw (x,0)--(x+bw,0)--(x+bw,y*u)--(x,y*u)--cycle withpen pencircle scaled 1.5pt withcolor c;
enddef;
numeric c[]; c1 := 2.0cm; c2 := 5.5cm; c3 := 9.0cm;
barra(c1-bw-0.05cm, 17.73, \MPcolor{darkblue});
barra(c2-bw-0.05cm, 134.16, \MPcolor{darkblue});
barra(c3-bw-0.05cm, 113.51, \MPcolor{darkblue});
barra(c1+0.05cm, 26.29, \MPcolor{redline});
barracortada(c2+0.05cm, 300, \MPcolor{redline});
barracortada(c3+0.05cm, 300, \MPcolor{redline});
label.top(textext("\tfx 17,7"), (c1-bw/2-0.05cm,17.73u));
label.top(textext("\tfx 134,2"), (c2-bw/2-0.05cm,134.16u));
label.top(textext("\tfx 113,5"), (c3-bw/2-0.05cm,113.51u));
label.top(textext("\tfx 26,3"), (c1+bw/2+0.05cm,26.29u));
label.bot(textext("\tfx Una velocidad"), (c1,-0.12cm));
label.bot(textext("\tfx Dos bandas"), (c2,-0.12cm));
label.bot(textext("\tfx Tres bandas"), (c3,-0.12cm));
label.bot(textext("Velocidades de la cuadrilla"), (w/2,-0.75cm));
label.lft(textext("Tiempo de la Etapa 2 (s)") rotated 90, (-1.1cm,h/2));
fill unitsquare scaled 0.3cm shifted (0.3cm,h+0.75cm) withcolor \MPcolor{darkblue};
label.rt(textext("\tfx 3 pickers"), (0.65cm,h+0.9cm));
fill unitsquare scaled 0.3cm shifted (2.9cm,h+0.75cm) withcolor \MPcolor{redline};
label.rt(textext("\tfx 4 pickers"), (3.25cm,h+0.9cm));
draw unitsquare scaled 0.3cm shifted (5.5cm,h+0.75cm) withpen pencircle scaled 1.5pt withcolor \MPcolor{redline};
label.rt(textext("\tfx 4 pickers, sin terminar en 300 s"), (5.85cm,h+0.9cm));
\stopMPcode}

% TODO: ojo acá--no es que no miramos la demora, aclarar eso. Si no, parece
% que solo testeamos el adelanto.
{\bf Pedidos por lote y peso de la finalización.} Se fijaron 26 pedidos por
lote. Durante la exploración, con un peso de finalización $W^{c} = 9$,
alrededor de dos tercios de los pedidos se completaban a tiempo, de modo que
en el objetivo actuaban tanto la penalidad por adelanto como la de demora.
Con ese peso, sin embargo, a los modelos lineales les convenía la espera
voluntaria (Sección~\in[sec:espera_voluntaria]), que las heurísticas no
pueden reproducir, y los objetivos de unos y otras dejaban de ser
comparables. Por eso, en la suite final el peso de finalización se fijó en
uno más que la cantidad de pedidos por lote ($W^{c} = 27$), de modo que
esperar nunca resulte conveniente. Con este peso los lotes terminan lo antes
posible y prácticamente todos los pedidos llegan a tiempo, por lo que en el
objetivo actúa casi solo el adelanto. El cambio también abarató la Etapa~2:
en los escenarios de grilla del perfil realista, su tiempo bajó de entre 86 y
203 segundos a entre 4 y 12 segundos.

{\bf Forma de la agenda.} Se fijaron 6 lotes para 3 {\em pickers}, de modo
que cada uno ejecute dos lotes y la secuenciación tenga efecto. Con las
velocidades en bandas, las formas más cargadas no terminaron en 25
minutos: con 6 lotes y 4 {\em pickers} la brecha al corte fue de 1,09\%,
con 7 lotes y 3 {\em pickers} de 2,6\%, y con 8 lotes y 3 {\em pickers}
de 9,48\%.

\stopsection

\startsection[title={¿Hasta qué tamaño de almacén podemos entregar resultados?}]

Con el perfil fijo, el único eje que queda libre es el tamaño del almacén.
Se buscó qué grillas llevan al modelo Split a resolver en un rango de 10 a
20 minutos, suficiente para exigirlo sin superar el tope de las pruebas.
Cada medición se tomó de a una, sin otras resoluciones en paralelo
(Figura~\in[fig:relev_calibracion]).

\placefigure[force][fig:relev_calibracion]{Tiempo de cada etapa del modelo Split en la calibración del perfil realista (6 lotes, 3 {\em pickers}, 26 pedidos por lote). La franja gris marca el rango buscado.}
{\startMPcode
numeric w, h, u, bw; w := 11cm; h := 6cm; u := h/26; bw := 1.6cm;
fill (0,10u)--(w,10u)--(w,20u)--(0,20u)--cycle withcolor \MPcolor{darkgray};
label.rt(textext("\tfx rango buscado: 10 a 20 min"), (0.15cm,17u));
drawarrow (0,0)--(w+0.4cm,0) withpen pencircle scaled 1pt;
drawarrow (0,0)--(0,h+0.4cm) withpen pencircle scaled 1pt;
for t = 0, 5, 10, 15, 20, 25:
  draw (0,t*u)--(-0.12cm,t*u) withpen pencircle scaled .5pt;
  label.lft(textext("\tfx " & decimal t), (-0.12cm,t*u));
endfor
def segmento(expr xa, ya, yb, cc) =
  fill (xa,ya*u)--(xa+bw,ya*u)--(xa+bw,yb*u)--(xa,yb*u)--cycle withcolor cc;
enddef;
numeric c[]; c1 := 2.0cm; c2 := 5.5cm; c3 := 9.0cm;
segmento(c1-bw/2, 0, 0.160, \MPcolor{darkblue});
segmento(c1-bw/2, 0.160, 3.564, \MPcolor{redline});
segmento(c2-bw/2, 0, 7.421, \MPcolor{darkblue});
segmento(c2-bw/2, 7.421, 11.818, \MPcolor{redline});
draw (c2-bw/2,7.421u)--(c2+bw/2,7.421u) withpen pencircle scaled 1pt withcolor white;
draw (c3-bw/2,0)--(c3+bw/2,0)--(c3+bw/2,25u)--(c3-bw/2,25u)--cycle withpen pencircle scaled 1.5pt withcolor \MPcolor{darkblue};
label.top(textext("\tfx 3,6 min"), (c1,3.564u));
label.top(textext("\tfx 11,8 min"), (c2,11.818u));
label.top(textext("\tfx más de 25 min"), (c3,25u));
label.bot(textext("\tfx 7×7"), (c1,-0.12cm));
label.bot(textext("\tfx 13×13"), (c2,-0.12cm));
label.bot(textext("\tfx 15×15"), (c3,-0.12cm));
label.bot(textext("Tamaño de la grilla"), (w/2,-0.75cm));
label.lft(textext("Tiempo (min)") rotated 90, (-1.0cm,h/2));
fill unitsquare scaled 0.3cm shifted (0.3cm,h+0.75cm) withcolor \MPcolor{darkblue};
label.rt(textext("\tfx Etapa 1 (ruteo)"), (0.65cm,h+0.9cm));
fill unitsquare scaled 0.3cm shifted (3.2cm,h+0.75cm) withcolor \MPcolor{redline};
label.rt(textext("\tfx Etapa 2 (asignación y secuencia)"), (3.55cm,h+0.9cm));
draw unitsquare scaled 0.3cm shifted (8.6cm,h+0.75cm) withpen pencircle scaled 1.5pt withcolor \MPcolor{darkblue};
label.rt(textext("\tfx Etapa 1 sin terminar"), (8.95cm,h+0.9cm));
\stopMPcode}

En la grilla de 7×7 casi todo el tiempo corresponde a la Etapa~2: el ruteo
tarda menos de 10 segundos. Al pasar a 13×13 la relación se invierte y la
Etapa~1 pasa a ser la más costosa, mientras que la Etapa~2 crece poco,
porque la forma de la agenda es la misma. En 15×15, la Etapa~1 no terminó
en 25 minutos. Estas mediciones definieron los escenarios de tamaño
creciente de las pruebas de estrés (Sección~\in[sec:escenarios_estres]).

\stopsection

\startsection[title={Alcance del modelo monolítico}, reference=sec:alcance_monolitico]

El modelo monolítico se evaluó sobre el mismo perfil, resolviendo cada
instancia también con el modelo Split para contar con su óptimo como
referencia. El monolítico solo logró probar el óptimo en instancias mucho
más chicas que las de la calibración (Figura~\in[fig:relev_frontera]): con
un tope de 10 minutos, hasta 3 lotes en una grilla de 6×6, 4 lotes en 5×5
y 6 lotes en 3×3. El modelo Split resolvió todos esos puntos en menos de
4 minutos. El objetivo de esta prueba de alcance es determinar cuál es la
frontera de factibilidad del modelo monolítico contra los del modelo Split sin 
ejecutar el monolítico en instancias más grandes.

\placefigure[force][fig:relev_frontera]{Resultado del modelo monolítico en cada punto del estudio, con un tope de 10 minutos (perfil realista). En azul, el tiempo que tardó en probar el óptimo; en rojo, la brecha al momento del corte.}
{\startMPcode
numeric cw, ch; cw := 1.9cm; ch := 0.85cm;
color azul, rojo, gris; azul := 0.25[white,\MPcolor{darkblue}]; rojo := 0.35[white,\MPcolor{redline}]; gris := \MPcolor{darkgray};
def celda(expr i, j, c, s) =
  fill ((j-1)*cw+1pt,(5-i)*ch+1pt)--(j*cw-1pt,(5-i)*ch+1pt)--(j*cw-1pt,(6-i)*ch-1pt)--((j-1)*cw+1pt,(6-i)*ch-1pt)--cycle withcolor c;
  draw textext("\tfx " & s) shifted ((j-0.5)*cw,(5.5-i)*ch);
enddef;
string fila[]; fila1 := "1 lote, 1 picker"; fila2 := "2 lotes, 1 picker"; fila3 := "3 lotes, 2 pickers"; fila4 := "4 lotes, 2 pickers"; fila5 := "6 lotes, 3 pickers";
for i = 1 upto 5: label.lft(textext("\tfx " & fila[i]), (-0.1cm,(5.5-i)*ch)); endfor
for j = 1 upto 5: label.top(textext("\tfx " & decimal (j+2) & "×" & decimal (j+2)), ((j-0.5)*cw,5ch)); endfor
celda(1,1,azul,"0,2 s"); celda(1,2,azul,"0,2 s"); celda(1,3,azul,"0,2 s"); celda(1,4,azul,"0,2 s"); celda(1,5,azul,"0,3 s");
celda(2,1,azul,"0,3 s"); celda(2,2,gris,"infactible"); celda(2,3,azul,"3,3 s"); celda(2,4,azul,"13,5 s"); celda(2,5,azul,"17,9 s");
celda(3,1,azul,"4,7 s"); celda(3,2,gris,"infactible"); celda(3,3,azul,"8,6 s"); celda(3,4,azul,"138,8 s"); celda(3,5,rojo,"11,4 \%");
celda(4,1,azul,"6,6 s"); celda(4,2,gris,"infactible"); celda(4,3,azul,"47,1 s"); celda(4,4,rojo,"13,2 \%"); celda(4,5,white,"--");
celda(5,1,azul,"192,1 s"); celda(5,2,gris,"infactible"); celda(5,3,rojo,"3,0 \%"); celda(5,4,white,"--"); celda(5,5,white,"--");
label.top(textext("Tamaño de la grilla"), (2.5cw,5ch+0.5cm));
fill unitsquare scaled 0.3cm shifted (0.5cm,-0.6cm) withcolor azul;
label.rt(textext("\tfx probó el óptimo (tiempo)"), (0.85cm,-0.45cm));
fill unitsquare scaled 0.3cm shifted (5.0cm,-0.6cm) withcolor rojo;
label.rt(textext("\tfx cortado a los 10 min (brecha)"), (5.35cm,-0.45cm));
fill unitsquare scaled 0.3cm shifted (0.5cm,-1.1cm) withcolor gris;
label.rt(textext("\tfx infactible"), (0.85cm,-0.95cm));
label.rt(textext("\tfx --: no se corrió"), (5.0cm,-0.95cm));
\stopMPcode}

Donde fue cortado, el monolítico ya había encontrado soluciones a menos del
0,3\% del óptimo. Su dificultad no está en encontrar buenas soluciones,
sino en demostrar que son óptimas. Con un tope de 40 minutos la situación
casi no cambió, y la única instancia adicional que cerró (6 lotes en 5×5)
le llevó 32 minutos, 8 veces lo que tardó el modelo Split. Las instancias
de 4×4 con más de un lote resultaron infactibles para ambos modelos, y las
celdas sin dato no se corrieron porque una grilla menor ya había alcanzado
el tope. Por estos resultados, en las pruebas de estrés la referencia
contra la que se comparan las heurísticas es el modelo Split.

\stopsection

\stopchapter
